\section{Related Literature}


Our framework relates to the literature on task-based models of technical change. 
Existing work generally views tasks as predetermined and technologically independent objects combined through a CES production function \citep{acemoglu2011skills, acemoglu2018, acemoglu2019automation, acemoglu2024tasks}.
In these settings, task-level comparative advantage pins down factor assignment and automation follows relative input efficiencies.
We depart from this literature by modeling production as a sequence of technologically interdependent steps, where AI chaining creates complementarities across adjacent steps that can overturn comparative advantage as the driver of who does what.

Our paper also relates to the literature on quantifying effects of new automation technologies on the labor market.
Earlier work develops measures of occupational exposure to AI \citep{frey2017future, brynjolfsson2018can, webb2020impact, felten2021occupational, eloundou2023gpts, tomlinson2025working}, typically by mapping technologies to specific tasks and then aggregating these task-level mappings linearly to construct occupation-level indices.
We depart from previous work by showing that linear aggregation obscures the crucial role of task interdependence.
In our framework, the gains from AI automation depend not only on how many steps in a job are exposed to AI, but also on whether those steps are clustered or dispersed within the production sequence.

We also contribute to the literature on friction-laden technology adoption and organizational complementarities.
When tasks exhibit complementarity, the return to improving one activity depends on performance elsewhere, in the spirit of O-ring theory of production \citep{kremer1993oring}.
In our setting AI chaining makes this logic salient, since the payoff to automating a step depends on whether neighboring steps can be bundled into the same AI block, creating threshold effects and discrete shifts in AI adoption and job design.
Our results align with \citet{gans2026oring}, who likewise emphasize task interdependence in O-ring production and show that complementarities can generate ``lumpy'' automation decisions via time/attention reallocation rather than workflow adjacency.

Our framework also provides a micro-foundation for the ``productivity J-curve,'' where adoption costs and transitional reorganization precede realized productivity gains \citep{brynjolfsson2021productivity}.
In our model, the gains from better AI are non-linear and arrive mainly once capability crosses thresholds that trigger reorganization and enable longer AI chains.
This dynamic is consistent with recent analyses of U.S. Census microdata, which show that AI adoption initially reduces productivity \citep{mcelheran2025rise}, as well as evidence from several other studies \citep{furman2019ai, tambe2012productivity, bonney2024tracking, mcelheran2024adoption}.

Finally, our work connects to the literature on the division of labor within a firm. 
Our approach builds on \cite{becker1992division}'s view that the division of labor is limited by coordination frictions and the difficulty of integrating specialized knowledge, and echoes \cite{dessein2006adaptive} in highlighting the tension between specialization and the need to coordinate interdependent tasks.\footnote{
\cite{ide2025ai} provides a related view on how AI in particular affects the structure of work. While we examine how it redefines job boundaries through a task-based framework, they embed AI as an autonomous problem solver in a knowledge hierarchy model \citep{garicano2000hierarchies, garicano2006organization} and study how it reshapes spans of control within the firm.
}
In our framework, AI can restructure the division of labor through chaining in two distinct ways.  
First, on the intensive margin, chaining can pull previously manual steps ``under the hood'' of a single AI-executed task, narrowing the span of activities requiring direct human attention and shifting responsibility within a role. 
Second, on the extensive margin, by lowering the cost of linking adjacent steps that lie across a job boundary, chaining can make it optimal to redraw that boundary, reassigning activities across roles and altering the pattern of realized hand-offs and coordination within the firm. 
Both channels effectively alter the skill requirements of workers and the expertise they must possess to perform their jobs \citep{autor2025expertise}.